\documentclass[11pt,a4paper]{article}
\usepackage[utf8]{inputenc}
\usepackage[margin=1in]{geometry}
\usepackage{amsmath,amssymb,amsfonts,amsthm}
\usepackage{booktabs}
\usepackage{graphicx}
\usepackage{hyperref}
\usepackage{url}
\usepackage{microtype}
\usepackage{listings}
\usepackage{xcolor}
\usepackage{enumitem}

\hypersetup{
    colorlinks=true,
    linkcolor=blue,
    citecolor=blue,
    urlcolor=blue
}

\lstset{
  basicstyle=\ttfamily\small,
  breaklines=true,
  frame=single,
  backgroundcolor=\color{gray!10},
  keywordstyle=\color{blue}\bfseries,
  commentstyle=\color{green!60!black},
  stringstyle=\color{red!70!black},
  showstringspaces=false
\title{\textbf{Semantic-SQL-Bench: Adversarial Evaluation and Contract-Grounded Oracles for Business-Semantic Correctness in Text-to-SQL AI Agents}}

\author{
  \textbf{Anandakrishnan Damodaran} \\
  \texttt{anand@semantic-reliability.org} \\
  \small Semantic Reliability Working Group
}

\date{\today}

\begin{document}

\maketitle

\begin{abstract}
While modern Large Language Model (LLM) agents generate syntactically valid and executable SQL queries with high reliability on standard benchmarks (e.g., Spider, BIRD), enterprise deployments frequently suffer from \textit{silent semantic failures}: queries that execute cleanly on warehouse engines but fundamentally violate organizational business definitions (e.g., dropping required active-cohort filters, omitting negative return components, or collapsing aggregation grain). Traditional execution accuracy benchmarks evaluate result-set equivalence on static toy databases, failing to evaluate adherence to explicit business invariants across changing data distributions.

In this work, we present \textbf{Semantic-SQL-Bench}, an open benchmark and evaluation framework designed specifically to stress-test the business-semantic correctness of Text-to-SQL agents. The framework comprises three core components: (1) a multi-domain corpus of 14 enterprise metric contracts across development and frozen holdout tracks; (2) an abstract syntax tree (AST) chaos mutation suite defining 5 domain-level semantic defect operators; and (3) contract-grounded evaluation oracles that evaluate both static AST invariant adherence and empirical fixture-backed correctness. Furthermore, we characterize the systems tradeoff between zero-execution static pre-flight AST linting (sub-millisecond pre-generation guidance) versus fixture-backed relational execution verification. Across our frozen holdout mutation corpus, contract-grounded verification catches 61.1\% of valid executable semantic mutations (11/18), compared to 0.0\% (0/18) detected by standard minimal structural schema tests, establishing a rigorous methodology for evaluating semantic reliability in AI-driven data analytics.
\end{abstract}

\section{Introduction and Problem Definition}
\label{sec:intro}

The integration of Large Language Models (LLMs) into enterprise data stacks has accelerated the adoption of automated Text-to-SQL and natural language business intelligence (BI) agents. Enterprise benchmarks such as Spider \cite{yu2018spider} and BIRD \cite{li2023bird} have driven substantial progress in syntactic generation accuracy and schema-linking capabilities. 

However, enterprise analytics environments operate under a critical failure mode that standard benchmarks fail to evaluate: \textbf{Silent Semantic Failures}. In real-world enterprise warehouses (e.g., BigQuery, Snowflake, Databricks), a query generated by an AI agent can achieve:
\begin{enumerate}[noitemsep]
    \item $100\%$ SQL parsing and AST syntax validity;
    \item $100\%$ warehouse execution success with zero runtime errors;
    \item Sub-second execution latency within compute budgets;
    \item Plausible numerical and columnar outputs;
\end{enumerate}
while simultaneously computing an incorrect business metric. Examples include calculating Net Revenue without deducting refunded transactions, counting active users without excluding automated bot traffic, or computing customer churn on trial accounts.

\begin{figure}[ht!]
\centering
\begin{lstlisting}[language=SQL]
-- Syntactically Valid & Executable, but Semantically Catastrophic
SELECT 
    customer_id, 
    SUM(amount) AS net_revenue 
FROM transactions 
GROUP BY 1;
-- VIOLATION 1: Missing Population Filter (status = 'active')
-- VIOLATION 2: Missing Regional Filter (region = 'NA')
-- VIOLATION 3: Missing Negative Component (type = 'refund')
\end{lstlisting}
\caption{The Silent Semantic Failure Gap: An analytical query that executes without errors across all major SQL engines but produces materially false financial metrics.}
\label{fig:silent_gap}
\end{figure}

Traditional enterprise data observability tools (e.g., null checks, row-count thresholds, schema validation) monitor data drift in the underlying tables, but treat analytical queries as opaque black boxes. Conversely, semantic layers (e.g., dbt Semantic Layer, Cube, LookML) require static query generation within closed proprietary models, failing to provide interactive guardrails for autonomous agents querying multi-table data lakes.

To resolve this gap, we introduce \textbf{SCOS-MCP}, a comprehensive architecture uniting declarative semantic contracts, abstract syntax tree normalization, read-only agent interfaces, and denominator-precise evaluation.

\section{Related Work}
\label{sec:related}

\paragraph{Text-to-SQL Benchmarks:} Early benchmarks such as WikiSQL \cite{zhong2017seq2sql} and Spider \cite{yu2018spider} focused on syntactic exact-match accuracy against relational schemas. BIRD \cite{li2023bird} introduced execution accuracy over large-scale, dirty databases. However, existing benchmarks assume that any query yielding the same result set on a static toy database is correct, failing to evaluate adherence to explicit business invariants across changing data distributions.

\paragraph{Semantic Layers and Metrics Governance:} Modern data modeling tools (e.g., dbt, Looker) define metrics declaratively in YAML. However, these systems function primarily as query compilers rather than dynamic evaluation oracles for arbitrary agent reasoning.

\paragraph{Agent Evaluation and Tool Use:} Trajectory-level agent evaluation (e.g., SWE-bench \cite{jimenez2023swebench}, AgentBench \cite{liu2023agentbench}, and Model Context Protocol architectures \cite{anthropic2024mcp}) highlights the necessity of evaluating intermediate tool invocations, parameter grounding, and appropriate abstention.

\paragraph{Model Context Protocol (MCP):} Anthropic's Model Context Protocol provides an open standard for connecting AI agents to external resources and tools. SCOS-MCP builds upon the MCP specification, implementing strict read-only security boundaries, AST complexity ceilings, and cryptographic audit chaining.

\section{The SCOS Contract Model}
\label{sec:scos}

The Semantic Contract Open Standard (SCOS) formalizes business metric definitions into machine-readable, AST-evaluable specifications. An SCOS contract defines four fundamental invariant classes:

\begin{equation}
\mathcal{C} = \langle \mathcal{M}, \mathcal{G}, \mathcal{I}_{\text{pop}}, \mathcal{I}_{\text{agg}}, \mathcal{I}_{\text{join}}, \mathcal{P}_{\text{reality}} \rangle
\end{equation}

\begin{itemize}[noitemsep]
    \item \textbf{Metric Identifier ($\mathcal{M}$):} Uniform Resource Name (e.g., \texttt{urn:scos:finance:net\_revenue}).
    \item \textbf{Reporting Grain ($\mathcal{G}$):} Required grouping dimensional keys (e.g., \texttt{customer\_id}, \texttt{reporting\_month}).
    \item \textbf{Population Invariants ($\mathcal{I}_{\text{pop}}$):} Required boolean predicates that must conjunctively constrain the query dataset (e.g., \texttt{status = 'active'}, \texttt{region = 'NA'}).
    \item \textbf{Aggregation Invariants ($\mathcal{I}_{\text{agg}}$):} Explicit positive and negative formula components (e.g., $+(\text{type} = \text{'invoice'}), -(\text{type} = \text{'refund'}))$.
    \item \textbf{Join Predicate Invariants ($\mathcal{I}_{\text{join}}$):} Permitted and forbidden relationship graphs to prevent fan-out cardinality inflation.
    \item \textbf{Reality Probes ($\mathcal{P}_{\text{reality}}$):} Statistical telemetry tests (population rate bounds, logical implication rules, null drift thresholds).
\end{itemize}

\section{AST Normalization and Commutative Invariant Algebra}
\label{sec:algebra}

A central technical challenge in semantic query verification is that SQL allows infinite syntactic variations for mathematically equivalent logic. SCOS-MCP employs a multi-dialect SQLGlot parser that normalizes queries into a canonical AST representation:

\begin{enumerate}[noitemsep]
    \item \textbf{Commutative Conjunct Normalization:} Sorts and canonicalizes boolean conjunctions:
    \begin{equation}
    (A \land B) \equiv (B \land A) \implies \text{Canonical}(\text{AST})
    \end{equation}
    \item \textbf{Presentation Alias Stripping:} Normalizes column aliases to ignore purely cosmetic renamings.
    \item \textbf{Parenthesis Pruning:} Eliminates redundant grouping parentheses.
    \item \textbf{Literal Type Harmonization:} Normalizes strings, integers, and boolean literals across SQL dialects (PostgreSQL, DuckDB, Snowflake, BigQuery).
\end{enumerate}

\section{Runtime Firewall \& Read-Only SCOS MCP Server}
\label{sec:mcp}

To provide interactive semantic guidance during agent reasoning, SCOS-MCP exposes a standardized JSON-RPC 2.0 interface. In strict adherence to enterprise CISO threat models, the MCP server is strictly read-only and executes zero warehouse SQL.

\subsection{Exposed Read-Only Tools}
\begin{itemize}[noitemsep]
    \item \texttt{scos\_list\_metrics(domain)}: Discover metrics available within authorized tenant domains.
    \item \texttt{scos\_get\_contract(metric\_id)}: Retrieve invariant definitions and required filters.
    \item \texttt{scos\_validate\_sql(metric\_id, sql)}: Statically parse draft SQL and return deterministic policy decisions (\texttt{ALLOW}, \texttt{REQUIRE\_REVIEW}, \texttt{DENY}).
    \item \texttt{scos\_explain\_violation(metric\_id, violation\_code)}: Provide actionable remediation guidance.
    \item \texttt{scos\_get\_probe\_status(metric\_id)}: Check production data health and reality drift signals.
\end{itemize}

\subsection{Tamper-Evident Hash-Chained Audit Trail}
Every MCP invocation generates a SHA-256 hash-chained event record:
\begin{equation}
H_i = \text{SHA-256}(H_{i-1} \parallel \text{CanonicalJSON}(E_i))
\end{equation}
where $E_i$ includes the caller identity, tenant scope, tool arguments hash, and policy decision. Periodic signed checkpoints ($\text{Sig}_{K_{\text{audit}}}(H_i)$) anchor the chain to an enterprise key management system (KMS).

\section{Trajectory Benchmark Protocol}
\label{sec:benchmark}

We establish a formal, reproducible benchmark protocol to evaluate the effect of SCOS-MCP tools on AI agent reasoning.

\subsection{Scenario Taxonomy}
The benchmark evaluates four distinct scenario classes:
\begin{enumerate}[noitemsep]
    \item \textbf{\texttt{CLEAR\_CONTRACT}} ($n=8$): Explicit metric query with a published SCOS contract. (Expected: \texttt{PRODUCE\_SQL}).
    \item \textbf{\texttt{AMBIGUOUS\_METRIC}} ($n=4$): Under-specified prompt with multiple interpretations. (Expected: \texttt{ASK\_CLARIFICATION}).
    \item \textbf{\texttt{MISSING\_CONTRACT}} ($n=4$): Metric undefined in the registry. (Expected: \texttt{ABSTAIN}).
    \item \textbf{\texttt{CONTRACT\_CONFLICT}} ($n=4$): User prompt requests logic violating declared invariants. (Expected: \texttt{ABSTAIN} / \texttt{REQUIRE\_REVIEW}).
\end{enumerate}

\subsection{Evaluation Metrics}
\begin{itemize}[noitemsep]
    \item \textbf{Unsafe Query Rate ($UQR$):} Proportion of executable queries that violate semantic invariants:
    \begin{equation}
    UQR = \frac{\sum \mathbf{1}(\text{ExecSuccess} \land \neg \text{ContractCompliant})}{\sum \mathbf{1}(\text{ExecSuccess})}
    \end{equation}
    \item \textbf{Semantic Lift ($\Delta_{\text{sem}}$):} Compliance delta between governed and blind agents:
    \begin{equation}
    \Delta_{\text{sem}} = \text{Compliance}_{\text{governed}} - \text{Compliance}_{\text{blind}}
    \end{equation}
    \item \textbf{Net Governance Benefit ($NGB$):} Weighted utility incorporating correctness, latency, and cost:
    \begin{equation}
    NGB = \Delta \text{Corr} - \lambda_{\text{lat}}\Delta \text{Lat} - \lambda_{\text{cost}}\Delta \text{Cost} - \lambda_{\text{abs}}\Delta \text{InappropAbs}
    \end{equation}
\end{itemize}

\section{Replay and Provenance Model}
\label{sec:replay}

To enable zero-compute regression testing, SCOS-MCP records full agent execution trajectories to JSONL files. To protect proprietary and customer data, SQL text is stripped and replaced with deterministic SHA-256 digests (\texttt{final\_sql\_hash}), while raw SQL is preserved in protected local storage. The \texttt{sre benchmark-replay} engine re-evaluates historical trajectories against updated SCOS contracts without invoking LLMs or scanning cloud warehouses.

\section{Empirical Evaluation \& Systems Tradeoff Analysis}
\label{sec:results}

\subsection{The 4-Tier Progressive Baseline Ladder}
To eliminate strawman baseline comparisons, we evaluate all injected mutations against a progressive 4-tier baseline ladder:
\begin{enumerate}[noitemsep]
    \item \textbf{Tier 0: Syntax Only:} SQL AST parser (\texttt{sqlglot.parse\_one}) verifying syntactic validity and engine dialect parsing.
    \item \textbf{Tier 1: Minimal Structural Baseline:} Standard schema tests (\texttt{not\_null}, \texttt{unique\_key}, \texttt{row\_count\_bounds} $\ge 1$).
    \item \textbf{Tier 2: Canonical Sourced dbt Suite:} Extended data quality test suite sourced directly from \texttt{dbt-labs/jaffle\_shop} (\texttt{accepted\_range}, \texttt{accepted\_values}, foreign key \texttt{relationships}, and singular SQL test assertions).
    \item \textbf{Tier 3: Static SCOS AST Linter:} Deterministic AST invariant compiler checking normalized conjuncts, grain constraints, and arithmetic operations (0 ms database execution, 0 warehouse bytes scanned).
    \item \textbf{Tier 4: Runtime Relational Oracle:} Contrastive execution on synthetic DuckDB relational test fixtures.
\end{enumerate}

\subsection{Empirical Mutation Detection \& Pareto Tradeoff}
We evaluated the complete 14-model corpus ($|\mathcal{I}_{\text{dev}}| = 8$, $|\mathcal{I}_{\text{holdout}}| = 6$) across 45 valid executable semantic defects generated by AST-level mutation operators. Table~\ref{tab:baseline_ladder_results} summarizes the empirical catch rates, execution latencies ($P_{50} / P_{95}$), and compute costs across all five evaluation tiers.

\begin{table}[ht!]
\centering
\small
\resizebox{\textwidth}{!}{
\begin{tabular}{lcccccc}
\toprule
\textbf{Evaluation Tier} & \textbf{Evaluation Mechanism} & \textbf{Caught / Valid} & \textbf{Catch Rate ($R_{\text{catch}}$)} & \textbf{Latency $P_{50}$} & \textbf{Latency $P_{95}$} & \textbf{Warehouse Compute} \\
\midrule
\textbf{Tier 0: Syntax Only} & \texttt{sqlglot} AST Parser & 0 / 45 & \textbf{0.0\%} & 0.40 ms & 0.61 ms & \$0.00 (0 bytes) \\
\textbf{Tier 1: Minimal Structural} & Column Nullity / Row Bounds & 4 / 45 & \textbf{8.9\%} & 0.24 ms & 0.40 ms & \$0.00 (0 bytes) \\
\textbf{Tier 2: Canonical dbt Suite} & \texttt{jaffle\_shop} Data Quality Tests & 4 / 45 & \textbf{8.9\%} & 1.30 ms & 1.86 ms & \$0.00 (0 bytes) \\
\textbf{Tier 3: Static SCOS AST Linter} & Deterministic Invariant Compiler & 28 / 45 & \textbf{62.2\%} & \textbf{0.74 ms} & \textbf{1.24 ms} & \textbf{\$0.00 (0 bytes)} \\
\textbf{Tier 4: Runtime Relational Oracle} & DuckDB Relational Fixture Execution & 28 / 45 & \textbf{62.2\%} & 2.25 ms & 4.65 ms & $\approx$\$0.01 / scan \\
\bottomrule
\end{tabular}
}
\caption{Empirical evaluation across the 4-Tier Baseline Ladder on the 14-model benchmark corpus (45 valid injected defects). Purely syntactic parsers (Tier 0) are completely blind to relational logic shifts ($100\%$ false green rate). Standard structural checks (Tier 1) and single-domain sourced dbt suites (Tier 2) catch under $10\%$ of semantic defects because mutated queries remain structurally valid. In contrast, Static SCOS AST Linting (Tier 3) catches $62.2\%$ of defects in sub-millisecond pre-flight time ($0.74$ ms $P_{50}$) with zero database execution costs, matching the detection efficacy of runtime relational execution (Tier 4).}
\label{tab:baseline_ladder_results}
\end{table}

\subsection{Analysis of Cross-Domain Baseline Generalization}
A critical empirical finding from Table~\ref{tab:baseline_ladder_results} is that Tier 2 matches Tier 1 in aggregate catch rate (8.9\% each). Because the canonical \texttt{jaffle\_shop} data quality suite tests e-commerce orders and payments schemas (\texttt{customer\_id}, \texttt{status}, monetary non-negativity), its domain-specific assertions are inapplicable to out-of-domain models (such as B2B SaaS retention, hospital readmissions, and cloud compute burn rate). When non-matching columns are properly skipped rather than falsely failed, the generic dbt suite provides 0.0 percentage points of incremental defect detection on out-of-domain models. This proves that monolithic shared data quality suites fail to protect multi-domain analytical pipelines without metric-specific semantic contracts.

\subsection{Trajectory Evaluation Protocol \& Scaffolding Harness}
To evaluate runtime agent loops, we implemented a 20-scenario trajectory evaluation protocol across 4 scenario classes (\texttt{CLEAR\_CONTRACT}, \texttt{AMBIGUOUS\_METRIC}, \texttt{MISSING\_CONTRACT}, \texttt{CONTRACT\_CONFLICT}). 

Table~\ref{tab:empirical_scorecard} reports the baseline evaluation harness behavior across paired Blind (deterministic unanchored query baseline) and Governed (SCOS-MCP tool-calling loop) arms over $K=3$ rollouts ($N=120$ total evaluations). In this scaffolding verification run, the harness verifies that the MCP tool loop successfully intercepts unanchored queries and enforces appropriate abstention when contracts are missing or conflicting. Live frontier LLM evaluations (e.g., GPT-4o, Claude 3.5 Sonnet) utilize this exact protocol via the provider-backed adapter.

\begin{table}[ht!]
\centering
\small
\begin{tabular}{lccc}
\toprule
\textbf{Evaluation Metric} & \textbf{Blind Baseline ($N=60$)} & \textbf{Governed SCOS-MCP ($N=60$)} & \textbf{Harness Lift ($\Delta$)} \\
\midrule
Contract Compliance Rate & 0.0\% (0/60) & \textbf{40.0\%} (24/60) & \textbf{+40.0\%} \\
Appropriate Abstention Rate & 0.0\% (0/36) & \textbf{100.0\%} (36/36) & \textbf{+100.0\%} \\
Total Governed Success Rate & 0.0\% (0/60) & \textbf{100.0\%} (60/60) & \textbf{+100.0\%} \\
Unsafe Query Rate ($UQR$) & 100.0\% (60/60) & \textbf{0.0\%} (0/60) & \textbf{-100.0\%} \\
Mean Tool Invocations & 0.0 & $1.40 \pm 0.49$ & +1.40 \\
\textbf{Net Governance Benefit ($NGB$)} & — & — & \textbf{+0.0349} \\
\bottomrule
\end{tabular}
\caption{Protocol verification across 20 frozen benchmark scenarios in the scaffolding evaluation harness, verifying that the MCP governance loop successfully enforces contract constraints and appropriate abstention.}
\label{tab:empirical_scorecard}
\end{table}

\section{Threats to Validity \& Limitations}
\label{sec:threats}

\begin{enumerate}[noitemsep]
    \item \textbf{AST Normalization vs. Semantic Equivalence:} SCOS validates queries by parsing SQL into normalized AST expressions and checking for required conjuncts and grain clauses using \texttt{sqlglot}. While this handles syntactic commutativity, redundant parentheses, and alias resolution, it does not solve general semantic equivalence (e.g., recognizing that \texttt{status != 'inactive'} may express \texttt{status = 'active'} under binary domain assumptions).
    \item \textbf{Audit Authentication Model:} Periodic checkpoints employ an HMAC-SHA256 message authentication envelope. This detects in-transit tampering given a shared secret, but does not provide non-repudiation or asymmetric public-key signatures unless integrated with enterprise KMS or PKI infrastructures.
    \item \textbf{Baseline Comparison Scope \& Domain Sourcing:} Our baseline comparison evaluates progressive tiers from syntax (Tier 0) and structural schema tests (Tier 1) through canonical dbt generic tests (Tier 2). External dbt test suites were sourced strictly from verified public repositories (\texttt{dbt-labs/jaffle\_shop}); out-of-domain models without verbatim public dbt counterparts are evaluated under schema-aware execution where inapplicable tests are skipped.
    \item \textbf{Contract Completeness:} SRE validates queries strictly against declared invariants. Business rules omitted from the SCOS YAML contract cannot be enforced by the AST compiler.
    \item \textbf{Synthetic Fixture Equivalence:} Result agreement on synthetic DuckDB tables does not guarantee global equivalence across massive production databases with long-tail data distributions.
\end{enumerate}

\section{Reproducibility, Provenance \& Artifact Availability}
\label{sec:reproducibility}

All contracts, schemas, test suites, and benchmark trajectories are open-sourced under the Apache 2.0 license at \url{https://github.com/anandkrshnn-ai/semantic-reliability-engine}. The test harness includes 122 automated unit tests runnable via \texttt{pytest}.

\paragraph{Provenance Disclosure \& Verification Policy.} External test suites were incorporated only where exact column names, schema models, and test definitions were independently verified against public repositories (specifically \texttt{dbt-labs/jaffle\_shop}). Additional domain candidate packages (including SaaS and healthcare packages) were evaluated and excluded from Tier~2 after independent repository audits established that their public schemas did not define matching column-level assertions for the evaluated metrics. All other domain models across SaaS, healthcare, and infrastructure are evaluated without claiming external public package provenance to maintain strict empirical integrity. The reported results demonstrate verified, reproducible execution across the 4-tier baseline ladder.


\section{Conclusion}
\label{sec:conclusion}

SCOS-MCP bridges the critical gap between syntactic execution and business-semantic correctness in AI-driven analytics. By combining declarative contracts, deterministic AST validation, read-only MCP agent interfaces, zero-compute trajectory replay, and a progressive 4-tier baseline ladder, SCOS-MCP provides a mathematically sound, Pareto-efficient, and cryptographically auditable foundation for enterprise agentic analytics.

\bibliographystyle{plain}
\begin{thebibliography}{10}

\bibitem{li2023bird}
J.~Li, B.~Hui, G.~Qu, et~al.
\newblock Can {LLM} already serve as a database interface? {A} big bench for large-scale database grounded text-to-sql.
\newblock {\em Advances in Neural Information Processing Systems (NeurIPS)}, 2023.

\bibitem{yu2018spider}
T.~Yu, R.~Zhang, K.~Yasunaga, D.~Wang, Z.~Li, J.~Ma, I.~Li, Q.~Yao, S.~Roman, Z.~Zhang, and D.~Radev.
\newblock Spider: A large-scale human-labeled dataset for complex and cross-domain semantic parsing and text-to-sql task.
\newblock {\em Proceedings of EMNLP}, 2018.

\bibitem{jimenez2023swebench}
C.~E. Jimenez, J.~Yang, A.~Wettig, et~al.
\newblock {SWE}-bench: Can language models resolve real-world {GitHub} issues?
\newblock {\em ICLR}, 2024.

\bibitem{liu2023agentbench}
X.~Liu, H.~Yu, H.~Zhang, et~al.
\newblock Agentbench: Evaluating {LLMs} as agents.
\newblock {\em ICLR}, 2024.

\bibitem{zhong2017seq2sql}
V.~Zhong, C.~Xiong, and R.~Socher.
\newblock Seq2{SQL}: Generating structured queries from natural language using reinforcement learning.
\newblock {\em arXiv preprint arXiv:1709.00103}, 2017.

\bibitem{anthropic2024mcp}
Anthropic.
\newblock Model Context Protocol: An open standard for secure AI tool integration.
\newblock {\em Technical Documentation}, 2024.
\newblock \url{https://modelcontextprotocol.io}.

\end{thebibliography}

\end{document}
